\begin{definition}
    Процесс $(W_t, t \geq 0)$ называется \textit{в$\acute{\text{и}}$неровским} (или процессом броуновского движения), если 
    \begin{enumerate}
        \item $W_0 = 0$ п.н.
        \item $W_t$ имеет независимые приращения
        \item $W_t - W_s \sim \mathcal{N}(0, t-s), \ \forall t > s \geq 0$
    \end{enumerate}
\end{definition}


\begin{statement}
Винеровский процесс (свойства 1-3) существует.
\end{statement}

\begin{proof}
По теореме о существовании процессов с независимыми приращениями достаточно проверить, что хар. функции правильно перемножаются:

Пусть $0 \leq s < u < t$, нужно проверить
$$
\varphi_{W_t - W_s}(\tau) \stackrel{?}{=} \varphi_{W_t - W_u}(\tau) \cdot \varphi_{W_u - W_s}(\tau)
$$
(Если $\xi \sim \Norm(0, \sigma^2)$, то её хар. функция $\varphi_\xi(t) = e^{-\frac{\sigma^2 t^2}{2}}$)

Так как $W_t - W_s \sim \mathcal{N}(0, t-s)$, то
$$
\varphi_{W_t - W_s}(\tau) = e^{- \frac{1}{2} \tau^2 (t-s)} = e^{- \frac{1}{2} \tau^2 (t-u)} \cdot e^{- \frac{1}{2} \tau^2 (u-s)} = \varphi_{W_t - W_u}(\tau) \cdot \varphi_{W_u - W_s}(\tau)
$$
\end{proof}

\begin{note}
    Нередко к требуемым свойствам винеровского процесса добавляют еще такое (доказательство в \ref{8.3}):
    \begin{equation*}
        \text{почти все траектории } w_{\bullet}(\omega) \text{ непрерывны}
    \end{equation*}
\end{note}

\begin{theorem}[эквивалентное определение винеровского процесса]
    Процесс $(W_t, t \geq 0)$ является винеровским $\Longleftrightarrow$ выполняются следующие свойства:
    \begin{enumerate}
        \item $W_t$~--- гаусовский процесс
        \item $\E W_t = 0, \ \forall t \geq 0$
        \item $cov(W_t, W_s) = min(t,s), \ \forall t, s \geq 0$
    \end{enumerate}
\end{theorem}

\begin{proof}
    Докажем пункт 3: при $t>s$ имеем

    \begin{gather*}
        t + s - 2 \E (W_tW_s) = \E (W_t - W_s)^2 = t - s \\
        \E (W_tW_s) = \min{(t, s)}
    \end{gather*}
\end{proof}
